% =====================================================================
%  v15 — first-draft snippets for sections currently MISSING or STUB
%  in v15_neurips_paper_v0.tex
%
%  Drop-in ready. Generated 2026-04-27.
%  Each block is preceded by its target location in the main paper.
% =====================================================================


% ---------------------------------------------------------------------
%  REPLACEMENT for §6.4 Broader impact (currently 1 short paragraph)
% ---------------------------------------------------------------------
\subsection{Broader impact}

\paragraph{Positive impact.}
A post-hoc, AUC-based diagnostic for preprocessing-induced flips
lowers the audit cost for biomedical ML pipelines, where many
published classifiers may be vulnerable to v5.1-style leakage.
DIAL has direct application in regulatory ML audits (FDA medical-
device re-validation, EMA cross-cohort review) and in ML-for-science
reproducibility programmes (NeurIPS Reproducibility Checklist, ICLR
ML+code reviews). The case study in \S\ref{sec:exp}.5 illustrates
that a leaky pipeline can produce statistically dramatic but
methodologically vacuous results; DIAL flags such cases at audit
time.

\paragraph{Negative / dual-use risk.}
DIAL is a diagnostic, not a corrective. We see no direct dual-use
risk: the method does not enable surveillance, manipulation of
clinical outcomes, or generation of harmful content. The closest
indirect risk is reviewer over-reliance on a single AUC-based
audit; we explicitly flag in \S\ref{sec:exp} (Limitations) that
DIAL is classifier- and CV-protocol-dependent, and that a low
DIAL value does not certify a pipeline as leak-free outside the
Theorem~\ref{thm:main} assumption regime.

\paragraph{Mitigations.}
We release both the metric implementation and the v5.2 leak-fix
audit script (\texttt{LinearComBat.fit}/\texttt{transform}
decomposition) with explicit fail-safe behaviour: under proper
LODO-fit-on-train-only, $\dial = 0$ at all $n$ for the leak-
inducing pipeline (Theorem~\ref{thm:main}). The v15 numerical
experiments are fully synthetic and reproducible from seeded
scripts; no human-subjects data are involved.

\paragraph{Scope of claim.}
DIAL fires under the Theorem~\ref{thm:main} conditions (Gaussian-
shared-$\Sigma$ class-conditionals, linear correction $T$).
Outside these assumptions, DIAL's behaviour is empirical and
modality-dependent; we benchmark four modalities in \S\ref{sec:exp}.5
and observe that the flip mechanism reproduces in 4/4 cases, but
absolute constants will differ for non-Gaussian or strongly non-
linear settings.


% ---------------------------------------------------------------------
%  NEW SECTION (insert before Appendix A): Ethics statement
% ---------------------------------------------------------------------
\section*{Ethics statement}

The v15 numerical experiments use no human-subjects data. The
v5.1 leak example referenced in \S\ref{sec:intro} and \S\ref{sec:exp}.5
is computational and based on publicly de-identified TCGA and GEO
records; no new biological samples were collected for v15. The
broader-impact discussion in \S6.4 covers dual-use considerations.
The author declares no conflict of interest.


% ---------------------------------------------------------------------
%  REPLACEMENT for "Code availability" (currently a stub)
% ---------------------------------------------------------------------
\section*{Code and data availability}
\label{app:code}

All code used to produce the v15 numerical experiments is at
\url{https://anonymous.4open.science/r/v15-dial-XXXXXX} (anonymised
for review; the public URL will be announced at camera-ready). The
five experiment scripts and their corresponding outputs are:

\begin{center}
\begin{tabular}{lll}
\toprule
Script & Section & Output checkpoint \\
\midrule
\texttt{v15\_theorem2\_verify.py} & \S\ref{sec:exp}.1 & \texttt{task1\_theorem2.json} \\
\texttt{v15\_stress\_test.py}     & \S\ref{sec:exp}.2 & \texttt{task2\_stress.json} \\
\texttt{v15\_tta\_benchmark.py}   & \S\ref{sec:exp}.3 & \texttt{task3\_tta.json} \\
\texttt{v15\_scaling.py}          & \S\ref{sec:exp}.4 & \texttt{task4\_scaling.json} \\
\texttt{v15\_cross\_domain.py}    & \S\ref{sec:exp}.5 & \texttt{task5\_cross\_domain.json} \\
\bottomrule
\end{tabular}
\end{center}

Each script writes a JSON checkpoint (recording exact wall time,
seed list, and headline metric) plus a long-form TSV that supports
the tabular numbers in the paper. No human-subjects data are used
in v15; the v5.1 leak example referenced in \S\ref{sec:intro} and
\S\ref{sec:exp}.5 is reproducible from public TCGA/GEO inputs but
is not required to validate any v15 claim.


% ---------------------------------------------------------------------
%  NEW APPENDIX SECTION: Hyperparameters table (closes Repro Q7)
% ---------------------------------------------------------------------
\section{Experimental hyperparameters}
\label{app:hparams}

\begin{center}
\begin{tabular}{llrl}
\toprule
Section & Hyperparameter & Value & Notes \\
\midrule
\S\ref{sec:exp}.1 & shift magnitudes        & 0,0.25,\dots,2.0 & 7 levels  \\
\S\ref{sec:exp}.1 & shift types             & 6                & none, cov, lab, par, perp, subspace  \\
\S\ref{sec:exp}.1 & repetitions             & 15               & per (type, mag) cell  \\
\S\ref{sec:exp}.1 & total runs              & 630              & 6$\times$7$\times$15  \\
\S\ref{sec:exp}.1 & wall                    & 3.4 s            & single-thread  \\
\midrule
\S\ref{sec:exp}.2 & feature dim $p$         & 64               &   \\
\S\ref{sec:exp}.2 & sample size $n$         & 300              & per cohort  \\
\S\ref{sec:exp}.2 & shift types             & 4                & cov, lab, concept, subspace  \\
\S\ref{sec:exp}.2 & magnitudes              & 5                & 0,0.25,0.5,0.75,1.0  \\
\S\ref{sec:exp}.2 & classifier families     & 5                & LogReg-$\ell_2$, RF, GB, XGB-substitute, MLP  \\
\S\ref{sec:exp}.2 & repetitions             & 20               & per cell  \\
\S\ref{sec:exp}.2 & total runs              & 2{,}000          &   \\
\S\ref{sec:exp}.2 & wall                    & 672 s            & 8 cores  \\
\midrule
\S\ref{sec:exp}.3 & feature dim $p$         & 64               &   \\
\S\ref{sec:exp}.3 & sample size $n$         & 300              & per cohort  \\
\S\ref{sec:exp}.3 & shift magnitude         & 0.8              & subspace-aligned  \\
\S\ref{sec:exp}.3 & methods                 & 10               & see Table in \S\ref{sec:exp}.3  \\
\S\ref{sec:exp}.3 & seeds                   & 5                &   \\
\midrule
\S\ref{sec:exp}.4 & encoder configurations  & 8                & 2K to 5.4M params  \\
\S\ref{sec:exp}.4 & shift magnitude         & 0.8              &   \\
\S\ref{sec:exp}.4 & seeds                   & 3                &   \\
\S\ref{sec:exp}.4 & wall                    & 96 s             &   \\
\midrule
\S\ref{sec:exp}.5 & domains                 & 4                & vision, NLP, biology, clinical  \\
\S\ref{sec:exp}.5 & dimensions per domain   & 32, 64, 128, 200 & shape-dependent  \\
\S\ref{sec:exp}.5 & seeds                   & 6                & per domain  \\
\bottomrule
\end{tabular}
\end{center}

Hardware: single workstation, AMD-class 8-core CPU, 32 GB RAM,
Linux 6.17. No GPU required. Total wall-clock for all five
experiments: $\approx 13$ min.


% ---------------------------------------------------------------------
%  REPRODUCIBILITY CHECKLIST (NeurIPS standard 14 questions)
% ---------------------------------------------------------------------
\section*{Reproducibility checklist}

\begin{enumerate}
\item \textbf{Claims and abstract accurately reflect contributions:}
  Yes. The abstract claims (R$^2 = 0.94$, ROC = 0.78, scaling
  negative result, 4/4 domains) are each backed by a checkpoint
  file referenced in \S\ref{sec:exp}.

\item \textbf{Limitations are clearly discussed:} Yes,
  \S\ref{sec:exp} (Limitations) explicitly lists the Gaussian
  assumption, classifier-dependence of DIAL, and the synthetic-only
  scope of v15 empirics.

\item \textbf{Theoretical claims have full proofs:} Theorem~\ref{thm:main}
  is proven in Appendix~A (\texttt{theorem2\_proof.tex}). \textbf{See
  Author Note: 10 TODO markers in the proof are flagged for
  verification.} A fallback restricted to Gaussian-only is
  available if any step does not go through.

\item \textbf{Did you describe assumptions and limits of theory:}
  Yes; Definition~1 in the proof setup states the Gaussian shared-
  $\Sigma$ assumption explicitly and \S\ref{sec:exp} (Limitations)
  acknowledges it.

\item \textbf{Experimental code provided:} Yes, five
  \texttt{v15\_*.py} scripts.

\item \textbf{Experimental data described:} Yes; all data are fully
  synthetic. Data generators are inline in each script.

\item \textbf{Hyperparameters specified:} Yes; see
  Appendix~\ref{app:hparams}.

\item \textbf{Random seeds reported:} Yes; checkpoints record
  \texttt{n\_seeds} per experiment (3 for scaling, 5 for TTA, 6 for
  cross-domain, 15 for theory validation, 20 for stress test).

\item \textbf{Variance reported:} Yes; \texttt{\_std} columns present
  in all TSV outputs.

\item \textbf{Compute resource described:} Yes; see
  Appendix~\ref{app:hparams} (single 8-core CPU, 32 GB, no GPU).
  Wall times in checkpoints (3 s to 672 s).

\item \textbf{Code license declared:} Yes; MIT (LICENSE in repo
  root).

\item \textbf{Data license declared:} Not applicable; data are
  synthetic and generated at runtime.

\item \textbf{Privacy / consent:} Not applicable; no human-subjects
  data.

\item \textbf{Theoretical fairness considered:} Yes; broader-impact
  statement (\S6.4) discusses dual-use risk, mitigations, and
  scope of claim.
\end{enumerate}


% ---------------------------------------------------------------------
%  CITATION INSTRUCTION FOR v17 (drop-in for any v17 paper that wants
%  to cite v15 DIAL methodology)
% ---------------------------------------------------------------------
%
%   "DIAL (Direction-Invariant AUC Leakage), a post-hoc diagnostic for
%   subspace-aligned conditional shift under linear batch correction
%   [Cook 2026 NeurIPS]; we apply DIAL to evaluate dark-matter
%   sub-cohort transfer to GSE27155 / GSE76039."
%
%   bibtex:
%   @inproceedings{cook2026dial,
%     title={DIAL: A post-hoc diagnostic for subspace-aligned
%       conditional shift under linear batch correction},
%     author={Cook, Seungho},
%     booktitle={NeurIPS 2026 [under review at v15-dial-XXXXXX]},
%     year={2026}
%   }
% ---------------------------------------------------------------------
